\section{Toward a quasi-polynomial hierarchy}
\label{sec:research}

\subsection{What this work removes from the complexity budget}
There are three different sizes that must not be conflated: the original
diagram size $n$, the encoded size of a current triangulated or uniform
handle state, and the expanded weight of a normal surface in that state.
The new normalization controls an operation in the second size while
avoiding dependence on the third as a unary parameter.  It does not yet
control how many states are discovered from the first.

In particular, a source with many interior vertices no longer forces an
immediate failure of the finite census method.  Theorem~\ref{thm:native}
removes all its top handles in polynomial encoded time, with no assumption
that their number is logarithmic.  The polynomial selection theorem also
removes a possible enumeration of dual trees in its stated column regime.
For a general matrix, however, the complete selection search retains its
$w^s$ parameter dependence.  The native geometric bound reduces its
original-column difficulty to cycle-compatible triple-unit columns, but
the augmented interface matrix can introduce additional difficulty.

The global composition issue can be stated without hiding an oracle.
Suppose a hierarchy discovery procedure makes at most $Q(n)$ state visits,
including revisits and rejected branches,
every state and emitted certificate has encoded size at most $S(n)$,
and the total discovery, transition-generation, local processing, and replay
work charged to each state visit costs at most $S(n)^d$, for a fixed
constant $d$.  A sequential implementation then costs at most
$Q(n)S(n)^{O(1)}$, plus the explicitly accounted input and output work.
If $Q(n)=2^{O((\log n)^3)}$ and $S(n)=n^{O(1)}$, this is the announced
quasi-polynomial scale.  This elementary composition observation is a
budgeting rule, not a new recognition theorem.

Lackenby's bound $L(g+1)^L$ gives the same scale if a suitable implemented
hierarchy can establish $L=O((\log n)^2)$, $g=n^{O(1)}$, and polynomial
encoded work at every node.  The linear hierarchy bound in the 2026
preprint, by itself, does not imply that logarithmic-square value of $L$
for arbitrary input.  Likewise a theorem that a small certificate exists
does not bound the work of finding it.  These are the principal gaps to
be closed, rather than details that can be suppressed in a local runtime.

\begin{table}[htbp]
\centering\small
\caption{Complexity obligations and their present status.}
\begin{tabular}{>{\raggedright\arraybackslash}p{0.30\textwidth}>{\raggedright\arraybackslash}p{0.60\textwidth}}
\toprule
Obligation & Status in this work\\
\midrule
Compressed cutting & Established in prior literature; the native pipeline
implements a source-certified form.\\
Remove native core three-handles & Proved and implemented in polynomial
encoded time for the admitted source contract.\\
Choose a tree for raw completion & Polynomial in the sparse active-column
case; complete parameterized search in general.\\
Transport recorded bundle ports & Proved and implemented for the current
exposed zero-/one-handle port charts.\\
Transport arbitrary hierarchy patterns & Open implementation and proof
obligation when patterns cross the cancellation support.\\
Find suitable surfaces efficiently & Not solved by this operator;
uniform discovery costs must be bounded.\\
Control state size and hierarchy depth & Must be established relative to
the original diagram, including all resets and certificates.\\
Decide every terminal state & The current raw kernel is sufficient but
incomplete; fourteen corpus cuts remain inconclusive.\\
\bottomrule
\end{tabular}
\end{table}

\subsection{Priority 1: remove the certified interface obstructions}
\paragraph{Question 1: which additional geometric moves discharge the three
saved annular stopping sets?}
The negative packets make this question sharply testable.  On the three
cuts in Table~\ref{tab:interfaces}, altering only the dual spanning tree
cannot work.  A next implementation should keep the authenticated source
and change a precisely named operation: the presentation of a bundle
attachment, a permitted disk-bundle absorption, or a certified lower-handle
move.  It should then show how the formerly nondeletable stopping rows
become removable or disappear under an actual geometric equivalence.
The success criterion is a new replayable genus census, not merely an
algebraic simplification of the row matrix.

\paragraph{Question 2: can an interface obstruction be localized to a
bounded subsurface?}
The saved stopping sets have sizes seven, four, and one.  Trace their
annular curves through the ordered boundary skeleton and determine the
smallest marked subsurface carrying them, together with its complementary
attachment data.  A useful theorem would show that a certified local
replacement has cost polynomial in the subsurface's encoded complexity
and preserves its gluing map relative to the complement.  A graph support
or homology calculation alone cannot guarantee this embedding statement.

\paragraph{Question 3: can primitive but nonunit attaching curves be
recognized with a compressed geometric certificate?}
The ball example shows why unsigned raw counts can stall on a primitive
curve.  One candidate extension is a sequence of boundary isotopies or
slides whose bands are specified by compressed paths and whose effect
on every other ordered attaching curve is independently replayed.
The target is not unrestricted free reduction of a word: cancelling
algebraic inverse letters can conceal geometric intersections with other
bands.  A successful theorem must bound both the move sequence and the
description of the transported curves.  The singleton annular obstruction
is a useful first native test case.

\subsection{Priority 2: solve the remaining native column geometry}
\paragraph{Question 4: is forest-compatible peeling polynomial for
cycle-compatible triple-unit columns?}
Propositions~\ref{prop:three} and~\ref{prop:cycle} show that the original
native columns lie in a much narrower class than general sparse matrices.
The remaining active shape is $(1,1,1)$, and its edges form a triangle,
a parallel pair plus a loop, or three loops.  The question is whether
these geometric restrictions allow an exact extension of the source-SCC
method.  They do not restore matroid exchange, as the explicit counterexample
shows.  An affirmative algorithm therefore needs additional structure;
a negative complexity result would require a reduction that respects
these constraints, preferably with a genuine native realization.

\paragraph{Question 5: is the number of triple-unit columns a useful
parameter?}
Let $k$ count the original columns of shape $(1,1,1)$ after the chosen
lower forest.  A strong target would be an algorithm of cost
$2^{\poly(k)}\poly(m,I)$, or an explicit smaller parameter based on the
interaction graph of those columns.  This would complement the current
$m^{O(s)}$ search when $s$ is large but $k$ is small.  The immediate
combinatorial model consists of ordinary reachability plus two-premise
activation rules.  Treating a two-premise rule as either of its individual
arcs is unsound; both prerequisites must be discharged, and the initial
deletion set must remain a graphic forest.

\paragraph{Question 6: can graphic rank be used earlier inside the
activation process?}
For ordinary source components, disjointness converts all obligations
into a partition matroid.  Ternary activation may create overlapping
obligations.  One approach is to maintain certified disjoint stopping
families and their contracted ranks throughout the search, extracting
short rank defects before branching.  Another is to identify separators
in the joint dual/dependency incidence graph and combine states recording
both connectivity and activation.  The research obligation is an explicit
state-count bound; declaring a separator small without bounding all
boundary partitions does not establish quasi-polynomial complexity.

\paragraph{Question 7: what exactly does changing the lower forest buy?}
The complete search theorem fixes the zero-/one-handle spanning forest.
Changing it changes which cocores survive, and can change the available
raw pivots.  Study paired choices of a lower forest and an upper dual
forest while retaining the source geometric cycles.  A first experiment
should exhaust both choices on the smallest authenticated failures,
record whether success depends on the lower choice, and save a replayable
certificate.  A theorem that bounds a representative family of lower
forests would be more valuable than unrestricted double enumeration.

\subsection{Priority 3: make hierarchy interfaces compositional}
\paragraph{Question 8: can arbitrary marked boundary patterns be transported
through the cancellation balls with bounded encoding growth?}
The present proof fixes exposed native ports.  In a general sutured or
patterned hierarchy, marked curves may cross the support of the selected
top cancellation.  The next geometric compiler must specify how such
curves move across the boundary disk $E$ and the replacement collar,
prove that the cyclic order and essential intersections are preserved,
and bound the growth of the compressed description.  A useful target is
a polynomial bound in the input pattern length and source chart size
per complete forest reduction, rather than a separate expansion at every
three-handle.

\paragraph{Question 9: can support-stratum reuse extend beyond the dual
tree without losing source authentication?}
Corollary~\ref{cor:support} identifies a valid reuse opportunity: the
dual graph and selected tree depend only on zero versus positive source
edge weights.  The full boundary words and bundle multiplicities need
not be constant on that stratum.  Seek a reusable symbolic template for
the retained chart whose variable lengths and endpoint offsets are
explicit arithmetic expressions in normal coordinates.  Each evaluation
must still authenticate the relevant source equalities.  This would
remove repeated structural work while keeping multiplicity-dependent
attachments honest.

\paragraph{Question 10: can local normalization be made confluent at a
useful level of abstraction?}
Different dual trees produce homeomorphic cores with potentially different
raw cancellation success.  Full geometric confluence is therefore too
strong for the present representation.  A more useful target is a canonical
summary of the remaining obstruction class, perhaps in terms of marked
boundary components and a certified family of stopping/rank witnesses.
The summary must preserve the information needed by subsequent cuts;
equal Euler characteristic, equal genus totals, or equal homology alone
does not suffice.

\subsection{Priority 4: close the discovery and global-size gaps}
\paragraph{Question 11: which supplied surfaces can be found with a proved
encoded-time bound?}
Start with a restricted source class where the existence theorem supplies
a normal disk or hierarchy surface of bounded combinatorial type.  Specify
the actual discovery algorithm, the bit lengths of its intermediate
coordinates, and the total unsuccessful work before a certificate is
found.  A normal-coordinate vector supplied by an experiment is an input
to this work's operator, not a discovery algorithm.  A promising bridge
to the hierarchy programme is a certified search over compressed surface
templates whose count can be bounded independently of expanded weight.

\paragraph{Question 12: can one charge every reset to a monotone global
quantity?}
The incoming reports already discuss causal reset budgets and cographic
persistence.  The next step should connect those budgets to the actual
geometric importer and this top normalization.  If an attachment change
invalidates a previous certificate or reintroduces branching, the reset
must have a reason recorded in a global potential.  Seek a statement
bounding the total number and cumulative encoded cost of resets along
every search path.  Independent local size bounds do not prevent an
algorithm from returning to exponentially many distinct states.

\paragraph{Question 13: can the total certificate, including all rejected
branches, be kept quasi-polynomial?}
Positive compressed certificates are often small while a search transcript
is large.  Distinguish a certificate needed to verify the eventual answer
from evidence needed to justify exhaustive failure of a branch.  The
rank defects here illustrate a beneficial compression of negative search.
Identify analogous short obstructions for other hierarchy operations,
and prove that their discovery cost is also controlled.  A concise
negative certificate discovered only after exponential search does not
resolve the runtime question.

\paragraph{Question 14: which terminal procedure should replace raw
cancellation when it is inadequate?}
The present kernel gives useful sufficient certificates, but the geometric
ball example proves its incompleteness in the broader handle category.
An eventual complete hierarchy needs a terminal method with a proved
worst-case budget for every permitted terminal state.  Candidate directions
include restricted geometric simplification, a certified comparison to
standard handlebodies, or a different hierarchy refinement.  The choice
should be evaluated against the actual terminal-state contract rather
than the full unrestricted handlebody problem.

\subsection{Priority 5: verification and credible performance measurement}
\paragraph{Question 15: what is the smallest independent trusted kernel?}
The current replay reduces dependence on discovery code but remains Python
and relies on the existing source geometry checker.  A serious next step
is a written specification for integer decoding, source incidence,
tree reduction, and ordered boundary transport, followed by a small
reference checker or proof-assistant formalization.  The parity constraint
is a useful independently derivable invariant, but should be checked on
uncapped source counts: capping three to two changes parity.

\paragraph{Question 16: where is the practical crossover on nonconstructed
inputs?}
The reported family isolates compression and the finite enumeration
family isolates tree choice.  Neither is a realistic input distribution
for all unknot recognition.  Extend measurement to a stratified native
corpus, reporting triangulation size, coordinate bit length, number of
interior vertices, number of triple-unit columns, interface size, genus
census, certificate size, wall time, and peak memory.  Compare complete
queries with randomized paired order and report regressions.  For an
end-to-end claim, include preprocessing, surface discovery, all rejected
cuts, final certificate production, and replay.

\subsection{Suggested next sequence of work}
The most informative immediate sequence is to resolve the three saved
interface obstructions with a named geometric extension; characterize
cycle-compatible ternary activation on native cores; and implement marked
pattern transport for the resulting reductions.  In parallel, one should
build an accounting layer that records discovery costs and state sizes
against the original diagram.  These tasks turn the current local theorem
and measured coverage improvement into reusable components of a hierarchy
argument.  The general quasi-polynomial recognition theorem remains the
goal to be proved after those obligations are met.
